\subsection{Směrnicový tvar rovnice přímky}
\label{subsec:Slope_intercept_line_equation}

V předchozích částech jsme přímku popisovali pomocí parametrického
vyjádření a obecné rovnice. Přímka však může být zároveň grafem funkce.
Z předchozího studia funkcí víme, že grafem lineární funkce
\[
    f(x)=ax+b
\]
je přímka.

Zapíšeme-li funkční předpis ve tvaru
\[
    y=kx+q,
\]
můžeme jej chápat nejen jako předpis funkce, ale také jako rovnici
přímky. Proměnná $x$ představuje první souřadnici bodu přímky
a hodnota $y$ jeho druhou souřadnici.

Tento způsob zápisu má jednu důležitou vlastnost. Každé hodnotě
$x\in\mathbb{R}$ je přiřazena právě jedna hodnota $y$. Přímka zapsaná
tímto způsobem proto nesmí být rovnoběžná s osou $y$, protože na svislé
přímce se může pro jednu hodnotu $x$ nacházet nekonečně mnoho bodů
s různými souřadnicemi $y$.

\begin{definition}\label{def:Slope_intercept_line_equation}
    Nechť $k,q\in\mathbb{R}$. Rovnici
    \begin{equation}\label{eq:Slope_intercept_line_equation}
        y=kx+q
    \end{equation}
    nazýváme \textbf{směrnicovým tvarem rovnice přímky}.

    Číslo $k$ nazýváme \textbf{směrnicí přímky} a číslo $q$
    \textbf{úsekem přímky na ose $y$}.
\end{definition}

Směrnicovým tvarem lze popsat každou přímku, která není rovnoběžná
s osou $y$. Pro takovou přímku jsou čísla $k$ a $q$ určena jednoznačně.
To znamená, že dva různé směrnicové tvary nemohou popisovat tutéž
přímku.

Svislé přímky mají rovnice tvaru
\[
    x=d,
    \qquad d\in\mathbb{R},
\]
a ve směrnicovém tvaru je zapsat nelze.

\subsubsection{Význam koeficientů $k$ a $q$}
\label{subsubsec:Meaning_slope_intercept_coefficients}

V rovnici
\[
    y=kx+q
\]
mají oba koeficienty přímý geometrický význam.

Dosadíme-li $x=0$, získáme
\[
    y=k\cdot0+q=q.
\]
Přímka proto vždy prochází bodem
\begin{equation}\label{eq:Y_intercept_point}
    Q[0;q].
\end{equation}

Číslo $q$ tedy určuje průsečík přímky s osou $y$. Změníme-li pouze
hodnotu $q$ a směrnici $k$ ponecháme stejnou, posuneme přímku ve
svislém směru, ale její směr se nezmění. Všechny přímky se stejnou
směrnicí jsou proto navzájem rovnoběžné nebo totožné.

Směrnice $k$ určuje, jak se změní souřadnice $y$, zvětšíme-li
souřadnici $x$ o jednu jednotku. Jestliže bod
\[
    A[x_A;y_A]
\]
leží na přímce, potom bod s první souřadnicí $x_A+1$ má druhou
souřadnici
\[
    y_A+k.
\]

Směrnice tedy vyjadřuje změnu druhé souřadnice připadající na jednotkovou
změnu první souřadnice.

\begin{figure}[H]
    \centering
    \begin{tikzpicture}[scale=1]

        % Osy
        \draw[<->] (-2,0) -- (5.5,0)
            node[right] {$x$};

        \draw[<->] (0,-0.8) -- (0,4.7)
            node[above] {$y$};

        \node[below left] at (0,0) {$O$};

        % Přímka y = kx + q
        \draw[thick, <->]
            (-1.5,0.1) -- (5,4)
            node[above right] {$p:y=kx+q$};

        % Body
        \coordinate (Q) at (0,1);
        \coordinate (R) at (2,1);
        \coordinate (A) at (2,2.2);

        \fill (Q) circle (2pt)
            node[left] {$Q[0;q]$};

        \fill (A) circle (2pt);

        % Pomocný pravoúhlý trojúhelník
        \draw[dashed] (Q) -- (R)
            node[midway, below] {$1$};

        \draw[dashed] (R) -- (A)
            node[midway, right] {$k$};

        % Úhel phi
        \draw (0.65,1)
            arc[start angle=0,end angle=31,radius=0.65];

        \node at (0.82,1.23) {$\varphi$};

        % Hodnota q na ose y
        \draw (-0.08,1) -- (0.08,1);
        \node[right] at (0,0.72) {$q$};

    \end{tikzpicture}

    \caption{Geometrický význam koeficientů $k$ a $q$
    ve směrnicovém tvaru přímky.}
    \label{fig:Slope_intercept_coefficients}
\end{figure}

Podle hodnoty směrnice mohou nastat následující případy:
\begin{itemize}
    \item je-li $k>0$, přímka je rostoucí;
    \item je-li $k<0$, přímka je klesající;
    \item je-li $k=0$, přímka je rovnoběžná s osou $x$.
\end{itemize}

\begin{remark}\label{rem:Special_slope_intercept_coefficients}
    Jestliže $q=0$, má přímka rovnici $y=kx$ a prochází počátkem soustavy souřadnic $O[0;0]$.

    Jestliže současně $k=0$, dostáváme rovnici
    \[
        y=0,
    \]
    která popisuje osu $x$.
\end{remark}

\subsubsection{Směrnice určená dvěma body}
\label{subsubsec:Slope_from_two_points}

Uvažujme dva různé body
\[
    A[x_A;y_A],
    \qquad
    B[x_B;y_B],
\]
které leží na přímce
\[
    y=kx+q.
\]

Pro jejich souřadnice platí
\[
    y_A=kx_A+q
\]
a
\[
    y_B=kx_B+q.
\]

Odečtením první rovnice od druhé získáme
\begin{align*}
    y_B-y_A &= kx_B+q-(kx_A+q),\\
    y_B-y_A &= kx_B-kx_A+q-q, \\
    y_B-y_A &= k(x_B-x_A).
\end{align*}

Jestliže $x_A\neq x_B$, můžeme rovnici vydělit rozdílem
$x_B-x_A$ a dostáváme
\begin{equation}\label{eq:Slope_from_two_points}
    k=
    \frac{y_B-y_A}{x_B-x_A}.
\end{equation}

Směrnici tedy můžeme chápat jako poměr změny souřadnice $y$ ke změně
souřadnice $x$:
\[
    k=\frac{\Delta y}{\Delta x}.
\]

Rozdíl
\[
    \Delta x=x_B-x_A
\]
představuje vodorovnou změnu mezi body, zatímco
\[
    \Delta y=y_B-y_A
\]
představuje změnu svislou.

Pokud platí $x_A=x_B$, je jmenovatel ve vztahu
\eqref{eq:Slope_from_two_points} nulový. Body $A$ a $B$ potom leží na
svislé přímce, jejíž směrnice není definována.

\begin{example}\label{exa:Slope_from_two_points}
    Určete směrnici přímky procházející body
    \[
        A[-2;1],
        \qquad
        B[4;4].
    \]

    Použijeme vztah \eqref{eq:Slope_from_two_points}:
    \[
        k
        =
        \frac{4-1}{4-(-2)}
        =
        \frac{3}{6}
        =
        \frac{1}{2}.
    \]

    Směrnice přímky je tedy
    \[
        k=\frac{1}{2}.
    \]

    To znamená, že při zvětšení souřadnice $x$ o dvě jednotky se
    souřadnice $y$ zvětší o jednu jednotku.
\end{example}

\subsubsection{Směrový úhel přímky a funkce tangens}
\label{subsubsec:Slope_angle_tangent}

Směrnici přímky můžeme vyjádřit také pomocí goniometrických funkcí.
Uvažujme přímku, která není rovnoběžná s osou $y$, a označme
$\varphi$ její \textbf{směrový úhel}, tedy úhel, který přímka svírá
s kladnou poloosou $x$.

Směrový úhel volíme z intervalu
\[
    \varphi\in\langle 0^\circ;180^\circ),
\]
přičemž pro nesvislou přímku platí
\[
    \varphi\neq90^\circ.
\]

Z pravoúhlého trojúhelníku na obrázku
\ref{fig:Slope_intercept_coefficients} získáme
\[
    \tan\varphi
    =
    \frac{\text{změna souřadnice }y}
         {\text{změna souřadnice }x}.
\]

Protože směrnice přímky je právě poměrem těchto změn, platí
\begin{equation}\label{eq:Slope_tangent_relation}
    k=\tan\varphi.
\end{equation}

Tento vztah propojuje analytickou geometrii s goniometrickými funkcemi.
Směrnici můžeme určit ze směrového úhlu a naopak směrový úhel můžeme
určit ze směrnice.

Je-li známá směrnice $k$, hledáme úhel splňující
\[
    \tan\varphi=k.
\]

Pro kladnou směrnici dostáváme ostrý úhel
\[
    0^\circ<\varphi<90^\circ.
\]

Pro zápornou směrnici musí být směrový úhel tupý:
\[
    90^\circ<\varphi<180^\circ.
\]

\begin{remark}\label{rem:Arctangent_negative_slope}
    Kalkulačka vrací při výpočtu $\arctan k$ obvykle hodnotu
    z intervalu
    \[
        (-90^\circ;90^\circ).
    \]
    Pro zápornou směrnici proto dostaneme záporný úhel. Chceme-li
    určit směrový úhel přímky z intervalu
    $\langle0^\circ;180^\circ)$, musíme k této hodnotě přičíst
    $180^\circ$.
\end{remark}

\begin{example}\label{exa:Direction_angle_from_slope}
    Určete směrový úhel přímky se směrnicí $k=-2$.

    Hledáme úhel $\varphi$, pro který platí
    \[
        \tan\varphi=-2.
    \]

    Kalkulačka dává
    \[
        \arctan(-2)\doteq-63^\circ26'.
    \]

    Protože je směrnice záporná, musí být směrový úhel přímky tupý.
    Přičteme tedy $180^\circ$:
    \[
        \varphi
        \doteq
        180^\circ-63^\circ26'
        =
        116^\circ34'.
    \]

    Směrový úhel přímky je přibližně
    \[
        \varphi\doteq116^\circ34'.
    \]
\end{example}

Pro vodorovnou přímku platí
\[
    k=0,
    \qquad
    \varphi=0^\circ.
\]

Svislá přímka má směrový úhel
\[
    \varphi=90^\circ,
\]
ale její směrnice není definována, protože funkce tangens není pro
úhel $90^\circ$ definována.

\subsubsection{Vztah směrnice a směrového vektoru}
\label{subsubsec:Slope_direction_vector}

Směrnice i směrový vektor určují směr přímky, a proto spolu musí
souviset.

Uvažujme přímku
\[
    y=kx+q.
\]
Zvětšíme-li souřadnici $x$ o jednu jednotku, změní se souřadnice $y$
o hodnotu $k$. Po přímce se tedy můžeme posunout o vektor
\begin{equation}\label{eq:Direction_vector_from_slope}
    \vec{u}=(1;k).
\end{equation}

Tento vektor je směrovým vektorem přímky. Stejně jako dříve jej můžeme
nahradit libovolným nenulovým násobkem:
\[
    \vec{u}=t(1;k),
    \qquad
    t\in\mathbb{R}\setminus\{0\}.
\]

Naopak má-li přímka směrový vektor
\[
    \vec{u}=(u_x;u_y),
\]
kde $u_x\neq0$, je její směrnice
\begin{equation}\label{eq:Slope_from_direction_vector}
    k=\frac{u_y}{u_x}.
\end{equation}

Tento vztah opět odpovídá poměru svislé a vodorovné změny.

\begin{example}\label{exa:Slope_from_direction_vector}
    Určete směrnici přímky, jejímž směrovým vektorem je
    \[
        \vec{u}=(-4;6).
    \]

    Podle vztahu \eqref{eq:Slope_from_direction_vector} platí
    \[
        k
        =
        \frac{6}{-4}
        =
        -\frac{3}{2}.
    \]

    Přímka má tedy směrnici
    \[
        k=-\frac{3}{2}.
    \]

    Stejnou směrnici bychom získali také z jednoduššího směrového
    vektoru
    \[
        \vec{v}=(2;-3),
    \]
    který je nenulovým násobkem původního vektoru.
\end{example}

Jestliže $u_x=0$, má směrový vektor tvar
\[
    \vec{u}=(0;u_y).
\]
Takový vektor určuje svislý směr. Přímka je tedy rovnoběžná s osou
$y$ a její směrnice není definována.

\subsubsection{Přímka určená bodem a směrnicí}
\label{subsubsec:Line_from_point_and_slope}

Známe-li směrnici přímky a jeden její bod, můžeme její rovnici snadno
určit.

Uvažujme přímku se směrnicí $k$, která prochází bodem
\[
    A[x_A;y_A].
\]

Směrnicový tvar má podobu
\[
    y=kx+q.
\]
Protože bod $A$ na přímce leží, musí jeho souřadnice rovnici splňovat:
\[
    y_A=kx_A+q.
\]

Odtud vyjádříme
\[
    q=y_A-kx_A.
\]

Dosazením do směrnicového tvaru dostáváme
\begin{align*}
    y
    &=
    kx+y_A-kx_A,\\
    y-y_A
    &=
    kx-kx_A.
\end{align*}

Po vytknutí směrnice $k$ získáme
\begin{equation}\label{eq:Point_slope_line_equation}
    y-y_A=k(x-x_A).
\end{equation}

\begin{definition}\label{def:Point_slope_line_equation}
    Rovnici
    \[
        y-y_A=k(x-x_A)
    \]
    nazýváme \textbf{bodově-směrnicovým tvarem rovnice přímky}.
    Popisuje přímku se směrnicí $k$, která prochází bodem
    $A[x_A;y_A]$.
\end{definition}

Výhodou tohoto zápisu je, že nemusíme nejprve určovat koeficient $q$.
Známý bod a směrnici dosadíme přímo do rovnice.

\begin{example}\label{exa:Line_from_point_and_slope}
    Určete směrnicový tvar přímky, která prochází bodem
    \[
        A[2;-1]
    \]
    a má směrnici
    \[
        k=\frac{3}{2}.
    \]

    Použijeme bodově-směrnicový tvar
    \[
        y-y_A=k(x-x_A).
    \]

    Po dosazení dostáváme
    \begin{align*}
        y-(-1)
        &=
        \frac{3}{2}(x-2),\\
        y+1
        &=
        \frac{3}{2}x-3,\\
        y
        &=
        \frac{3}{2}x-4.
    \end{align*}

    Hledaná přímka má směrnicový tvar
    \[
        p:y=\frac{3}{2}x-4.
    \]

    Správnost ověříme dosazením bodu $A[2;-1]$:
    \begin{align*}
        -1 &= \frac{3}{2} \cdot 2 - 4, \\
        -1 &= 3 - 4, \\
        -1 &= -1.  
    \end{align*}
\end{example}

\begin{remark}\label{rem:All_lines_through_point_slope}
    Všechny nesvislé přímky procházející bodem $A[x_A;y_A]$
    můžeme zapsat jako $y-y_A=k(x-x_A)$, $k\in\mathbb{R}$.
    Tato soustava však neobsahuje svislou přímku procházející bodem
    $A$. Tu musíme přidat samostatně: $x=x_A$.
\end{remark}

\subsubsection{Přímka určená dvěma body}
\label{subsubsec:Slope_intercept_line_from_two_points}

Jsou-li dány dva různé body
\[
    A[x_A;y_A],
    \qquad
    B[x_B;y_B],
\]
nejprve pomocí vztahu \eqref{eq:Slope_from_two_points} určíme směrnici
přímky. Poté použijeme jeden ze zadaných bodů v bodově-směrnicovém
tvaru.

Tento postup lze použít pouze tehdy, když
\[
    x_A\neq x_B.
\]

\begin{example}\label{exa:Slope_intercept_line_from_two_points}
    Určete směrnicový tvar přímky, která prochází body $A[-2;1], B[4;4]$.

    Směrnici jsme již určili v příkladu
    \ref{exa:Slope_from_two_points}:
    \[
        k=\frac{1}{2}.
    \]

    Použijeme bod $A[-2;1]$:
    \begin{align*}
        y-1
        &=
        \frac{1}{2}(x-(-2)),\\
        y-1
        &=
        \frac{1}{2}(x+2),\\
        y-1
        &=
        \frac{1}{2}x+1,\\
        y
        &=
        \frac{1}{2}x+2.
    \end{align*}

    Hledaná přímka má rovnici
    \[
        p:y=\frac{1}{2}x+2.
    \]

    Dosazením bodu $B[4;4]$ ověříme
    \[
        4
        =
        \frac{1}{2}\cdot4+2
        =
        2+2
        =
        4.
    \]
\end{example}

\subsubsection{Příslušnost bodu k přímce}
\label{subsubsec:Point_on_slope_intercept_line}

Bod
\[
    P[x_P;y_P]
\]
leží na přímce
\[
    p:y=kx+q
\]
právě tehdy, když jeho souřadnice splňují rovnici
\begin{equation}\label{eq:Point_on_slope_intercept_line}
    y_P=kx_P+q.
\end{equation}

\begin{example}\label{exa:Points_on_slope_intercept_line}
    Rozhodněte, zda body $C[2;3]$ a $D[-2;0]$ leží na přímce $p:y=\frac{1}{2}x+2$.

    Pro bod $C[2;3]$ po dosazení $x$-ové souřadnice dostáváme
    \[
        \frac{1}{2}\cdot2+2
        =
        1+2
        =
        3.
    \]
    Vypočtená hodnota odpovídá $y$-nové souřadnici bodu $C$, a proto
    \[
        C\in p.
    \]

    A pro bod $D[-2;0]$ dostáváme
    \[
        \frac{1}{2}\cdot(-2)+2
        =
        -1+2
        =
        1.
    \]
    Druhá souřadnice bodu $D$ je však rovna $0$, a proto
    \[
        D\notin p.
    \]
\end{example}

\subsubsection{Průsečíky se souřadnicovými osami}
\label{subsubsec:Slope_intercept_axis_intersections}

Průsečík přímky
\[
    y=kx+q
\]
s osou $y$ můžeme určit přímo z její rovnice. Protože na ose $y$ platí
$x=0$, dostáváme bod
\begin{equation}\label{eq:Slope_intercept_y_axis_intersection}
    P_y[0;q].
\end{equation}

Chceme-li určit průsečík s osou $x$, dosadíme $y=0$:
\[
    0=kx+q.
\]

Je-li $k\neq0$, získáme
\[
    x=-\frac{q}{k}.
\]
Průsečík s osou $x$ má tedy souřadnice
\begin{equation}\label{eq:Slope_intercept_x_axis_intersection}
    P_x\left[-\frac{q}{k};0\right].
\end{equation}

\begin{example}\label{exa:Slope_intercept_axis_intersections}
    Určete průsečíky přímky $p:y=-2x+4$ se souřadnicovými osami.

    Průsečík s osou $y$ získáme přímo z hodnoty $q=4$:
    \[
        P_y[0;4].
    \]

    Pro průsečík s osou $x$ položíme $y=0$:
    \begin{align*}
        0&=-2x+4,\\
        2x&=4,\\
        x&=2.
    \end{align*}

    Průsečík s osou $x$ je
    \[
        P_x[2;0].
    \]
\end{example}